l\documentclass[12pt,letterpaper]{article}
\usepackage{graphicx,textcomp}
\usepackage{natbib}
\usepackage{setspace}
\usepackage{fullpage}
\usepackage{color}
\usepackage[reqno]{amsmath}
\usepackage{amsthm}
\usepackage{fancyvrb}
\usepackage{amssymb,enumerate}
\usepackage[all]{xy}
\usepackage{endnotes}
\usepackage{lscape}
\newtheorem{com}{Comment}
\usepackage{float}
\usepackage{hyperref}
\newtheorem{lem} {Lemma}
\newtheorem{prop}{Proposition}
\newtheorem{thm}{Theorem}
\newtheorem{defn}{Definition}
\newtheorem{cor}{Corollary}
\newtheorem{obs}{Observation}
\usepackage[compact]{titlesec}
\usepackage{dcolumn}
\usepackage{tikz}
\usetikzlibrary{arrows}
\usepackage{multirow}
\usepackage{xcolor}
\newcolumntype{.}{D{.}{.}{-1}}
\newcolumntype{d}[1]{D{.}{.}{#1}}
\definecolor{light-gray}{gray}{0.65}
\usepackage{url}
\usepackage{listings}
\usepackage{color}

\definecolor{codegreen}{rgb}{0,0.6,0}
\definecolor{codegray}{rgb}{0.5,0.5,0.5}
\definecolor{codepurple}{rgb}{0.58,0,0.82}
\definecolor{backcolour}{rgb}{0.95,0.95,0.92}

\lstdefinestyle{mystyle}{
	backgroundcolor=\color{backcolour},   
	commentstyle=\color{codegreen},
	keywordstyle=\color{magenta},
	numberstyle=\tiny\color{codegray},
	stringstyle=\color{codepurple},
	basicstyle=\footnotesize,
	breakatwhitespace=false,         
	breaklines=true,                 
	captionpos=b,                    
	keepspaces=true,                 
	numbers=left,                    
	numbersep=5pt,                  
	showspaces=false,                
	showstringspaces=false,
	showtabs=false,                  
	tabsize=2
}
\lstset{style=mystyle}
\newcommand{\Sref}[1]{Section~\ref{#1}}
\newtheorem{hyp}{Hypothesis}

\title{Problem Set 3}
\date{Due: March 25, 2026}
\author{Applied Stats II}


\begin{document}
	\maketitle
	\section*{Instructions}
	\begin{itemize}
	\item Please show your work! You may lose points by simply writing in the answer. If the problem requires you to execute commands in \texttt{R}, please include the code you used to get your answers. Please also include the \texttt{.R} file that contains your code. If you are not sure if work needs to be shown for a particular problem, please ask.
\item Your homework should be submitted electronically on GitHub in \texttt{.pdf} form.
\item This problem set is due before 23:59 on Wednesday March 25, 2026. No late assignments will be accepted.
	\end{itemize}

	\vspace{.25cm}
\section*{Question 1}
\vspace{.25cm}
\noindent We are interested in how governments' management of public resources impacts economic prosperity. Our data come from \href{https://www.researchgate.net/profile/Adam_Przeworski/publication/240357392_Classifying_Political_Regimes/links/0deec532194849aefa000000/Classifying-Political-Regimes.pdf}{Alvarez, Cheibub, Limongi, and Przeworski (1996)} and is labelled \texttt{gdpChange.csv} on GitHub. The dataset covers 135 countries observed between 1950 or the year of independence or the first year forwhich data on economic growth are available ("entry year"), and 1990 or the last year for which data on economic growth are available ("exit year"). The unit of analysis is a particular country during a particular year, for a total $>$ 3,500 observations. 

\begin{itemize}
	\item
	Response variable: 
	\begin{itemize}
		\item \texttt{GDPWdiff}: Difference in GDP between year $t$ and $t-1$. Possible categories include: "positive", "negative", or "no change"
	\end{itemize}
	\item
	Explanatory variables: 
	\begin{itemize}
		\item
		\texttt{REG}: 1=Democracy; 0=Non-Democracy
		\item
		\texttt{OIL}: 1=if the average ratio of fuel exports to total exports in 1984-86 exceeded 50\%; 0= otherwise
	\end{itemize}
	
\end{itemize}
\newpage
\noindent Please answer the following questions:

\begin{enumerate}
	\item Construct and interpret an unordered multinomial logit with \texttt{GDPWdiff} as the output and "no change" as the reference category, including the estimated cutoff points and coefficients.
	I wrangled the data to set the reference categories and then ran the model with multinom. 
	\lstinputlisting[language=R, firstline=41, lastline=51]{PS3_Answers_IB.R}
	
	Results: 
	Call:
	multinom(formula = GDPWdiffcat ~ REG + OIL, data = gdp_data)
	
	Coefficients:
						(Intercept)      REG      OIL
	Decrease    3.805370 1.379282 4.783968
	Increase    4.533759 1.769007 4.576321
	
	Std. Errors:
						(Intercept)       REG      OIL
	Decrease   0.2706832 0.7686958 6.885366
	Increase   0.2692006 0.7670366 6.885097
	
	Residual Deviance: 4678.77 
	AIC: 4690.77 
	
	On average the log odds of a decrease relative to the odds of no change all other variables being zero is 3.80, while the log odds of an increase relative to the odds of no change is 4.53.
	
	On average, all other variables held constant, being a democracy increases the log odds of a decrease in GDP by about 1.38 relative to the chance of no change, and a log odds of an increase in GDP by about 1.77 relative to no change. If the average ratio of fuel exports to total exports in 1984-86 exceeded half, there is a 4.7 increase in the log odds of GDP decreasing relative to the chance of no change, and a 4.5 increase in log odds of GDP increasing relative to the chance of no change.
	
	\item Construct and interpret an ordered multinomial logit with \texttt{GDPWdiff} as the outcome variable, including the estimated cutoff points and coefficients.
	\lstinputlisting[language=R, firstline=53, lastline=54]{PS3_Answers_IB.R}
	Results: 
	
	Coefficients:
				Value Std. Error t value
	REG  0.4102    0.07518   5.456
	OIL -0.1788    0.11546  -1.549
	
	Intercepts:
											Value    Std. Error t value 
	No change|Decrease  -5.3199   0.2523   -21.0865
	Decrease|Increase   -0.7036   0.0476   -14.7932
	
	Residual Deviance: 4686.606 
	AIC: 4694.606 
	
The intercepts mean that the odds of no change relative to a decrease are -5.3199, while the log odds of a decrease relative to an increase are -0.7036.
	
All else being equal, on average being a democracy increases the log odds of GDP difference increasing by .41, while having the OIL variable be true decreases the log odds by .1788.
	
\end{enumerate}

\section*{Question 2} 
\vspace{.25cm}

\noindent Consider the data set \texttt{MexicoMuniData.csv}, which includes municipal-level information from Mexico. The outcome of interest is the number of times the winning PAN presidential candidate in 2006 (\texttt{PAN.visits.06}) visited a district leading up to the 2009 federal elections, which is a count. Our main predictor of interest is whether the district was highly contested, or whether it was not (the PAN or their opponents have electoral security) in the previous federal elections during 2000 (\texttt{competitive.district}), which is binary (1=close/swing district, 0="safe seat"). We also include \texttt{marginality.06} (a measure of poverty) and \texttt{PAN.governor.06} (a dummy for whether the state has a PAN-affiliated governor) as additional control variables. 

\begin{enumerate}
	\item [(a)]
	Run a Poisson regression because the outcome is a count variable. Is there evidence that PAN presidential candidates visit swing districts more? Provide a test statistic and p-value.
	\lstinputlisting[language=R, firstline=65, lastline=68]{PS3_Answers_IB.R}
	Coefficients:
										Estimate Std. Error z value Pr(>|z|)    
	(Intercept)         	 		-3.81023    0.22209 -17.156   <2e-16 ***
	competitive.district	  -0.08135    0.17069  -0.477   0.6336    
	marginality.06       		-2.08014    0.11734 -17.728   <2e-16 ***
	PAN.governor.06    		-0.31158    0.16673  -1.869   0.0617 .  
	---
	Signif. codes:  0 ‘***’ 0.001 ‘**’ 0.01 ‘*’ 0.05 ‘.’ 0.1 ‘ ’ 1
	
	(Dispersion parameter for poisson family taken to be 1)
	
	Null deviance: 1473.87  on 2406  degrees of freedom
	Residual deviance:  991.25  on 2403  degrees of freedom
	AIC: 1299.2
	
	Number of Fisher Scoring iterations: 7
	
	The test stat and p values for the intercept are -17.156 and  <2e-16, -0.477 and 0.6336  for competitive.district, -17.728 and <2e-16 for marginality, and -1.869 and 0.0617 for pan.governor.
	

	\item [(b)]
	Interpret the \texttt{marginality.06} and \texttt{PAN.governor.06} coefficients.
	For each unit of marginality increase, there is a -2.08014 log percentage decrease in the number of visits. If the state has a PAN-affiliated governor, there is a -.31158 log percentage decrease in the number of visits.
	
	\item [(c)]
	Provide the estimated mean number of visits from the winning PAN presidential candidate for a hypothetical district that was competitive (\texttt{competitive.district}=1), had an average poverty level (\texttt{marginality.06} = 0), and a PAN governor (\texttt{PAN.governor.06}=1).
	
	E[number of visits] = exp(-3.81023 + -0.08135 + -2.08014)
	Estimated number of visits are 0.002549852, which would be 0. 
\end{enumerate}

\end{document}
